\section{Experimental setup}
\label{sec:experiments}

\subsection{Two studies}

The experiments answer two different questions, and they disagree. Keeping
them separate is what makes the disagreement legible rather than
contradictory.

\begin{description}[leftmargin=0em,itemsep=0.35em]
  \item[Study G (general capability).] Train on the 279-task mixture,
    evaluate on seven tasks never trained on. \emph{Does it work on a
    schema you have no labels for?} This is the harder claim and the
    reference system's main advantage.
  \item[Study S (specific use case).] Train on 395 examples of the
    healthcare router, evaluate on its 450-item test split. \emph{Given
    labels for your actual problem, how good does it get?}
\end{description}

\subsection{Backbones}

\begin{center}
\begin{tabular}{lll}
\toprule
 & encoder & decoder \\
\midrule
model & \texttt{ModernBERT-base} & \texttt{Qwen3-1.7B} \\
parameters & 150M & 1{,}725M \\
attention & bidirectional & causal + block-diagonal \\
readout & linear at marker & $\text{logit}(\texttt{yes}) - \text{logit}(\texttt{no})$ \\
\bottomrule
\end{tabular}
\end{center}

``Decoder'' names the attention pattern, not the usage: nothing is
generated token by token, and inference is a single forward pass with a
parallel read at marker positions.

\subsection{Training configurations}

Study G uses one epoch over the mixture, batch 16 (encoder) or 8
(decoder), context 2048, label-space augmentation at $\pi{=}0.5$ with cap
$M{=}120$. Study S uses six epochs over 395 examples. Four adaptation
regimes are compared in \S\ref{sec:res-router}: encoder full fine-tune,
decoder full fine-tune, decoder head-only, and decoder LoRA ($r{=}16$,
$\alpha{=}32$, on all attention and MLP projections, $\text{lr}=2\times
10^{-4}$).

\subsection{Metrics}

Accuracy, macro-F1, and multi-label exact-set match, with bootstrap
confidence intervals over examples. For comparability across tasks with
very different $K$ we report \emph{multiple of chance}, $\text{acc} \cdot
K$, since raw accuracy is not comparable: $0.28$ on a 77-way menu is
strong and on a 4-way menu is worse than guessing.

Two metric choices are load-bearing.

\paragraph{AUROC on binary tasks.} Argmax accuracy on a \noul{} describes
the $0.5$ threshold as much as the model. A task can rank almost perfectly
and still score exactly the class balance by putting all probability mass
on one side. We report both, and the gap between them is a calibration
result rather than a transfer result.

\paragraph{Chance is not always the right floor.} A skewed task can be
beaten by always predicting its most common label. Where that baseline
exceeds $1/K$ we test against it explicitly, which changes one of our
conclusions (\S\ref{sec:res-heldout}).

\paragraph{Paired comparison.} Against the reference system we use exact
McNemar on the same items, reporting $b_{10}$ (we are right, they are not)
and $b_{01}$. An aggregate delta cannot distinguish a real difference from
two systems disagreeing equally in both directions.

\subsection{Reproducibility}

All runs are on one H100. The contamination guard runs and logs on every
training run. Checkpoints record their backbone, adapter rank, prompt
preamble, option template, and fitted temperatures, because the same
weights scored under a different prompt are a different system. Artifacts
are listed in Appendix~\ref{app:artifacts}.
